\documentclass[aps,prb,onecolumn,nofootinbib,superscriptaddress]{revtex4-2}

\usepackage{amsmath,amssymb,amsthm,mathtools,bm}
\usepackage{booktabs}
\usepackage[colorlinks=true,linkcolor=blue,citecolor=blue,urlcolor=blue]{hyperref}
\usepackage{microtype}

\newcommand{\I}{\mathbf{1}}
\newcommand{\Tr}{\operatorname{Tr}}
\newcommand{\Ran}{\operatorname{Ran}}
\newcommand{\ket}[1]{\lvert #1\rangle}
\newcommand{\bra}[1]{\langle #1\rvert}
\newcommand{\eps}{\varepsilon}
\newcommand{\cH}{\mathcal{H}}
\newcommand{\cK}{\mathcal{K}}

\theoremstyle{plain}
\newtheorem{proposition}{Proposition}

\theoremstyle{definition}
\newtheorem{remark}{Remark}

\begin{document}

\title{Regularized Choi-Covariance Contractions for the Class-A Adaptive Fermionic Circuit}
\author{Asadullah Bhuiyan}
\date{\today}

\begin{abstract}x
We formalize the regulated Gaussian contraction calculation underlying a
cumulative Choi-covariance representation of realized trajectories in the
class-A adaptive fermionic circuit.  A local rank-one Wannier-mode operation is
represented by a two-leg covariance tensor
\(
K(\eta_1,\eta_2)
=
\left(\begin{smallmatrix}
\eta_1P & \I-P\\
\I-P & \eta_2P
\end{smallmatrix}\right)
\),
where \(P=\ket{\chi}\bra{\chi}\) and
\(\eta_1,\eta_2\in\{\pm1\}\) select the four measurement or correction
branches.  Direct composition with an accumulated Choi covariance is
apparently obstructed because \(K_{LL}=\eta_1P\) is singular on the
orthogonal complement of the measured mode.  Introducing
\(K_{LL}^{(\eps)}=\eta_1P+\eps(\I-P)\), performing the Schur-complement
contraction, and subsequently taking \(\eps\to0\) produces a finite update
whenever
\(
\Tr(P\Sigma_{RR})-\eta_1\neq0
\).
We give the blockwise calculation, correct a dimensional ambiguity in the
lower-right block of the original handwritten derivation, and state the
result entirely in basis-independent projector notation.  This cumulative
operator construction is distinct from the support-reduced state-covariance
dynamics implemented by the repository's canonical Markov-circuit engine.
\end{abstract}

\maketitle
\tableofcontents

\section{Circuit-Operator Viewpoint}

Consider one realized sequence of local measurement and feedback events.  Its
unnormalized Gaussian circuit operator is denoted by \(\hat C_t\), so that the
normalized physical state at time step \(t\) is
\begin{equation}
  \ket{\psi_t}
  =
  \frac{\hat C_t\ket{\psi_0}}
       {\|\hat C_t\ket{\psi_0}\|},
  \qquad
  \hat C_{t+1}=\hat K_{t+1}\hat C_t .
  \label{eq:operator-trajectory}
\end{equation}
The object considered below is a two-leg Gaussian Choi covariance for
\(\hat C_t\),
\begin{equation}
  \Sigma_t
  =
  \begin{pmatrix}
    \Sigma_{LL} & \Sigma_{LR}\\
    \Sigma_{RL} & \Sigma_{RR}
  \end{pmatrix},
  \qquad
  \Sigma_{RL}=\Sigma_{LR}^{\dagger}.
  \label{eq:Sigma-blocks}
\end{equation}
The \(L\) and \(R\) labels refer to the input and output legs of the
cumulative operator tensor, rather than to spatial regions.

This viewpoint must be distinguished from the state-covariance formulation
used for production dynamics in this repository.  There the dynamical variable
is the top-layer covariance
\begin{equation}
  G_{ij}=2\langle c_i^\dagger c_j\rangle-\delta_{ij},
  \label{eq:G-convention}
\end{equation}
which is propagated directly by
\texttt{classA\_U1FGTN.run\_markov\_circuit(...)} or its GPU analogue.
The present calculation instead asks how a full realized operator
\(\hat C_t\) would be updated and retained as a Choi tensor.

\section{Local Adaptive Branches}

Let \(\ket{\chi}\) be a normalized single-particle Wannier mode and let
\begin{equation}
  P=\ket{\chi}\bra{\chi},
  \qquad
  Q=\I-P.
  \label{eq:PQ}
\end{equation}
Thus \(P\) is rank one and \(Q\) projects onto all modes orthogonal to the
locally interrogated channel.  The four local Gaussian branches written in
the handwritten calculation can be assembled into
\begin{equation}
  K(\eta_1,\eta_2)
  =
  \begin{pmatrix}
    K_{LL} & K_{LR}\\
    K_{RL} & K_{RR}
  \end{pmatrix}
  =
  \begin{pmatrix}
    \eta_1P & Q\\
    Q & \eta_2P
  \end{pmatrix},
  \qquad
  \eta_1,\eta_2\in\{+1,-1\}.
  \label{eq:local-K}
\end{equation}
In the sign convention of the handwritten notes, the four branch assignments
encoded by Eq.~\eqref{eq:local-K} are listed in Table~\ref{tab:branches}.
\begin{table}[b]
  \caption{\label{tab:branches}
  Rank-one local operations encoded by Eq.~\eqref{eq:local-K}.  The
  \(+\) and \(-\) labels on the mode indicate upper- and lower-band
  Observable-Wannier channels, not the signs \(\eta_{1,2}\).}
  \begin{ruledtabular}
  \begin{tabular}{ccc}
  Operation & Physical role & \((\eta_1,\eta_2)\)\\
  \colrule
  \(\hat N_-\) &
  lower-band mode already occupied & \((-1,+1)\)\\
  \(\hat\chi_-^\dagger(1-\hat N_-)\) &
  fill a lower-band vacancy & \((+1,+1)\)\\
  \(1-\hat N_+\) &
  upper-band mode already empty & \((+1,-1)\)\\
  \(\hat\chi_+\hat N_+\) &
  remove an upper-band particle & \((-1,-1)\)
  \end{tabular}
  \end{ruledtabular}
\end{table}

In the present implementation of the circuit, each site supplies the
Observable-Wannier channels \(\chi_{A+}\), \(\chi_{A-}\), \(\chi_{B+}\),
and \(\chi_{B-}\).  The desired state empties
the upper-band modes \(A+\) and \(B+\), while filling the lower-band modes
\(A-\) and \(B-\).  Consequently the table is the operator-level
counterpart of the measurement and feedback logic implemented for
\texttt{Ap}, \texttt{Am}, \texttt{Bp}, and \texttt{Bm}.

\section{Two-Leg Composition Formula}

Contract the right leg of the accumulated tensor \(\Sigma\) with the left
leg of a new local tensor \(K\).  The standard Gaussian two-leg contraction
gives
\begin{align}
  \Sigma'
  &=
  K\circ\Sigma
  \nonumber\\
  &=
  \begin{pmatrix}
    \Sigma_{LL} & 0\\
    0 & \eta_2P
  \end{pmatrix}
  -
  \begin{pmatrix}
    \Sigma_{LR} & 0\\
    0 & Q
  \end{pmatrix}
  \begin{pmatrix}
    \Sigma_{RR} & -\I\\
    -\I & \eta_1P
  \end{pmatrix}^{-1}
  \begin{pmatrix}
    \Sigma_{RL} & 0\\
    0 & Q
  \end{pmatrix}.
  \label{eq:singular-composition}
\end{align}
The expression is formally singular because the internal lower-right block
\(\eta_1P\) has rank one.  Inverting that block would require an inverse of
\(P\) on \(\Ran Q\), where it vanishes identically.

To define the contraction as a limit, replace the singular block by
\begin{equation}
  A_\eps
  =
  \eta_1P+\eps Q,
  \qquad
  A_\eps^{-1}
  =
  \eta_1P+\frac{1}{\eps}Q,
  \qquad
  \eps\neq0.
  \label{eq:regulator}
\end{equation}
Let
\begin{equation}
  R\equiv \Sigma_{RR},
  \qquad
  \cK_\eps
  =
  \begin{pmatrix}
    R & -\I\\
    -\I & A_\eps
  \end{pmatrix}.
  \label{eq:internal-regulated}
\end{equation}
For fixed nonzero \(\eps\), block inversion through the Schur complement
\begin{equation}
  S_\eps=R-A_\eps^{-1}
  \label{eq:S-eps}
\end{equation}
yields
\begin{equation}
  \cK_\eps^{-1}
  =
  \begin{pmatrix}
  S_\eps^{-1}
  &
  S_\eps^{-1}A_\eps^{-1}
  \\[2pt]
  A_\eps^{-1}S_\eps^{-1}
  &
  A_\eps^{-1}
  +
  A_\eps^{-1}S_\eps^{-1}A_\eps^{-1}
  \end{pmatrix}.
  \label{eq:internal-inverse}
\end{equation}
Substitution into Eq.~\eqref{eq:singular-composition} reduces the problem to
four regulated limits.

\section{Calculation in a Projector-Adapted Basis}

Choose a temporary computational basis adapted to
\begin{equation}
  \cH=\Ran P\oplus\Ran Q,
  \qquad
  \dim\Ran P=1.
  \label{eq:PQ-splitting}
\end{equation}
In this basis,
\begin{equation}
  P=
  \begin{pmatrix}
    1&0\\0&0
  \end{pmatrix},
  \qquad
  Q=
  \begin{pmatrix}
    0&0\\0&\I_Q
  \end{pmatrix},
  \qquad
  R=
  \begin{pmatrix}
    a&b^\dagger\\
    b&C
  \end{pmatrix},
  \label{eq:R-adapted}
\end{equation}
where
\begin{equation}
  a\in\mathbb{C},
  \qquad
  b\in\mathbb{C}^{(N-1)\times 1},
  \qquad
  C\in\mathbb{C}^{(N-1)\times(N-1)}.
  \label{eq:reduced-dimensions}
\end{equation}
For Hermitian \(R\), \(a\) is real and \(C=C^\dagger\).  Define
\begin{equation}
  d=a-\eta_1.
  \label{eq:d-adapted}
\end{equation}
The following discussion assumes \(d\neq0\).

Using Eq.~\eqref{eq:regulator},
\begin{equation}
  S_\eps
  =
  \begin{pmatrix}
    d & b^\dagger\\
    b & C-\eps^{-1}\I_Q
  \end{pmatrix}.
  \label{eq:S-adapted}
\end{equation}
Since
\begin{equation}
  (C-\eps^{-1}\I_Q)^{-1}
  =
  -\eps(\I_Q-\eps C)^{-1}
  =
  -\eps\I_Q-\eps^2C+O(\eps^3),
  \label{eq:Q-inverse-expansion}
\end{equation}
the block inverse of \(S_\eps\) gives the finite combinations needed by the
Choi contraction:
\begin{align}
  S_\eps^{-1}
  &\longrightarrow
  \begin{pmatrix}
    d^{-1}&0\\
    0&0
  \end{pmatrix},
  \label{eq:limit-one}
  \\
  S_\eps^{-1}A_\eps^{-1}Q
  &\longrightarrow
  \begin{pmatrix}
    0&d^{-1}b^\dagger\\
    0&-\I_Q
  \end{pmatrix},
  \label{eq:limit-two}
  \\
  QA_\eps^{-1}S_\eps^{-1}
  &\longrightarrow
  \begin{pmatrix}
    0&0\\
    d^{-1}b&-\I_Q
  \end{pmatrix},
  \label{eq:limit-three}
  \\
  Q\!\left(
    A_\eps^{-1}
    A_\eps^{-1}S_\eps^{-1}A_\eps^{-1}
  \right)Q
  &\longrightarrow
  \begin{pmatrix}
    0&0\\
    0&d^{-1}bb^\dagger-C
  \end{pmatrix}.
  \label{eq:limit-four}
\end{align}
Although the inverse \(A_\eps^{-1}\) diverges on \(\Ran Q\), its divergent
contribution cancels in the combinations selected by the external \(Q\)
factors of the local tensor.

\subsection{Block Updates in the Adapted Basis}

Equations~\eqref{eq:limit-one}--\eqref{eq:limit-four} immediately imply
\begin{align}
  \Sigma_{LL}'
  &=
  \Sigma_{LL}
  -
  \Sigma_{LR}
  \begin{pmatrix}
    d^{-1}&0\\0&0
  \end{pmatrix}
  \Sigma_{RL},
  \label{eq:LL-adapted}
  \\
  \Sigma_{LR}'
  &=
  \Sigma_{LR}
  \begin{pmatrix}
    0&-d^{-1}b^\dagger\\
    0&\I_Q
  \end{pmatrix},
  \label{eq:LR-adapted}
  \\
  \Sigma_{RL}'
  &=
  \begin{pmatrix}
    0&0\\
    -d^{-1}b&\I_Q
  \end{pmatrix}
  \Sigma_{RL},
  \label{eq:RL-adapted}
  \\
  \Sigma_{RR}'
  &=
  \boxed{
  \begin{pmatrix}
    \eta_2&0_{1\times(N-1)}\\[4pt]
    0_{(N-1)\times1}
    &
    C-\dfrac{bb^\dagger}{a-\eta_1}
  \end{pmatrix}
  }.
  \label{eq:RR-adapted}
\end{align}

\section{Dimensional Correction to the Lower-Right Block}

The handwritten calculation identifies the correct reduced \(Q\)-sector
Schur complement,
\begin{equation}
  C-\frac{bb^\dagger}{a-\eta_1}.
  \label{eq:reduced-RR}
\end{equation}
There is, however, an important dimensional qualification.  The matrices
\(C\) and \(bb^\dagger\) in Eq.~\eqref{eq:reduced-RR} act on
\(\Ran Q\) and hence have dimensions \((N-1)\times(N-1)\).  By contrast,
\(P\), \(Q\), and \(\Sigma_{RR}'\) act on the full \(N\)-dimensional
single-particle space.  Thus an expression that adds \(\eta_2P\) directly
to Eq.~\eqref{eq:reduced-RR}, or multiplies Eq.~\eqref{eq:reduced-RR} by the
full-space projector \(Q\) without declaring an embedding, is dimensionally
incomplete.

The corrected statement is Eq.~\eqref{eq:RR-adapted}: the scalar
\(\eta_2\) occupies the \(P\)-sector, while the reduced Schur complement is
embedded as the \(Q\)-sector block.  In physical terms, the selected mode is
reset to the covariance prescribed by \(\eta_2\), and the modes not selected
by \(P\) inherit the conditional Schur-complement covariance induced by their
previous correlations with that mode.

\section{Basis-Independent Update}

The adapted basis serves only to calculate the limit.  Its final result can
be written entirely using \(P\), \(Q\), and the original Choi blocks.  First,
the reduced quantities appearing above obey
\begin{align}
  a
  &=
  \Tr(PR)
  =
  \bra{\chi}R\ket{\chi},
  \label{eq:a-invariant}
  \\
  b
  &\longleftrightarrow
  QR\ket{\chi},
  \label{eq:b-invariant}
  \\
  bb^\dagger
  &\longleftrightarrow
  QRPRQ,
  \label{eq:bb-invariant}
  \\
  C
  &\longleftrightarrow
  QRQ\big|_{\Ran Q}.
  \label{eq:C-invariant}
\end{align}
The arrows in Eqs.~\eqref{eq:b-invariant}--\eqref{eq:C-invariant} indicate
that the left-hand objects are matrices in a basis of \(\Ran Q\), whereas
the right-hand expressions are their full-space operator embeddings.

\begin{proposition}[Regulated rank-one Choi contraction]
\label{prop:reg-contraction}
Let \(P=\ket{\chi}\bra{\chi}\) be rank one, \(Q=\I-P\), and
\[
  K(\eta_1,\eta_2)
  =
  \begin{pmatrix}
    \eta_1P&Q\\
    Q&\eta_2P
  \end{pmatrix}.
\]
Let \(\Sigma\) be as in Eq.~\eqref{eq:Sigma-blocks}, write
\(R=\Sigma_{RR}\), and suppose
\begin{equation}
  d
  =
  \Tr(PR)-\eta_1
  =
  \bra{\chi}R\ket{\chi}-\eta_1
  \neq 0.
  \label{eq:d-invariant}
\end{equation}
Then the \(\eps\to0\) limit of the contraction obtained from
\(K_{LL}^{(\eps)}=\eta_1P+\eps Q\) exists and is
\begin{equation}
  \boxed{
  \Sigma'
  =
  \begin{pmatrix}
  \displaystyle
  \Sigma_{LL}
  -
  \frac{\Sigma_{LR}P\Sigma_{RL}}{d}
  &
  \displaystyle
  \Sigma_{LR}Q
  -
  \frac{\Sigma_{LR}PRQ}{d}
  \\[15pt]
  \displaystyle
  Q\Sigma_{RL}
  -
  \frac{QRP\Sigma_{RL}}{d}
  &
  \displaystyle
  \eta_2P+QRQ-\frac{QRPRQ}{d}
  \end{pmatrix}.
  }
  \label{eq:full-invariant-update}
\end{equation}
\end{proposition}

For clarity, the four operator-valued blocks of
Eq.~\eqref{eq:full-invariant-update} are
\begin{align}
  \boxed{
  \Sigma_{LL}'
  =
  \Sigma_{LL}
  -
  \frac{\Sigma_{LR}P\Sigma_{RL}}
       {\Tr(P\Sigma_{RR})-\eta_1}
  },
  \label{eq:LL-invariant}
  \\[5pt]
  \boxed{
  \Sigma_{LR}'
  =
  \Sigma_{LR}Q
  -
  \frac{\Sigma_{LR}P\Sigma_{RR}Q}
       {\Tr(P\Sigma_{RR})-\eta_1}
  },
  \label{eq:LR-invariant}
  \\[5pt]
  \boxed{
  \Sigma_{RL}'
  =
  Q\Sigma_{RL}
  -
  \frac{Q\Sigma_{RR}P\Sigma_{RL}}
       {\Tr(P\Sigma_{RR})-\eta_1}
  },
  \label{eq:RL-invariant}
  \\[5pt]
  \boxed{
  \Sigma_{RR}'
  =
  \eta_2P
  +
  Q\Sigma_{RR}Q
  -
  \frac{Q\Sigma_{RR}P\Sigma_{RR}Q}
       {\Tr(P\Sigma_{RR})-\eta_1}
  }.
  \label{eq:RR-invariant}
\end{align}
Every term in Eqs.~\eqref{eq:LL-invariant}--\eqref{eq:RR-invariant} is an
operator on the full \(N\)-dimensional single-particle space.  No
dimension-dependent insertion maps or implied reduced-space products remain.

\subsection{Sherman--Morrison Form of the \texorpdfstring{$RR$}{RR} Block}
\label{subsec:SM-RR}

The \(Q\)-sector of \(\Sigma_{RR}'\) admits a compact inverse
representation.  Whenever \(R=\Sigma_{RR}\) is invertible, the
Sherman--Morrison formula applied to the rank-one perturbation
\(-\eta_1 P\) of \(R^{-1}\) gives
\begin{equation}
  \bigl(R^{-1}-\eta_1 P\bigr)^{-1}
  =
  R
  +
  \frac{\eta_1\,RPR}
       {1-\eta_1\langle\chi\rvert R\lvert\chi\rangle}.
  \label{eq:SM-R-inverse}
\end{equation}
Because \(\langle\chi\rvert R\lvert\chi\rangle = d+\eta_1\) and
\(\eta_1^2=1\), the denominator collapses:
\begin{equation}
  1-\eta_1(d+\eta_1)
  =
  1-\eta_1 d-\eta_1^2
  =
  -\eta_1 d,
  \label{eq:SM-denom}
\end{equation}
so Eq.~\eqref{eq:SM-R-inverse} becomes
\begin{equation}
  \bigl(R^{-1}-\eta_1 P\bigr)^{-1}
  =
  R-\frac{RPR}{d}.
  \label{eq:SM-simplified}
\end{equation}
Projecting onto \(\Ran Q\) on both sides,
\begin{equation}
  Q\bigl(R^{-1}-\eta_1 P\bigr)^{-1}Q
  =
  QRQ-\frac{QRPRQ}{d},
  \label{eq:SM-projected}
\end{equation}
which is exactly the \(Q\)-sector of \(\Sigma_{RR}'\) in
Eq.~\eqref{eq:RR-invariant}.  The full updated right-right block
therefore reads
\begin{equation}
  \boxed{
  \Sigma_{RR}'
  =
  \eta_2 P
  +
  Q\bigl(\Sigma_{RR}^{-1}-\eta_1 P\bigr)^{-1}Q.
  }
  \label{eq:RR-SM-form}
\end{equation}
Equation~\eqref{eq:RR-SM-form} makes the singular-branch condition
transparent: the Sherman--Morrison denominator for inverting
\(\Sigma_{RR}^{-1}-\eta_1 P\) is \(-\eta_1 d\), which vanishes if and only
if \(d=0\).  The regulated contraction is therefore well-defined on precisely
those branches for which the inverse in
Eq.~\eqref{eq:RR-SM-form} exists.

\subsection{Regularized Resolvent Form of All Blocks}
\label{subsec:regularized-resolvent-blocks}

The inverse form in Eq.~\eqref{eq:RR-SM-form} can be written without first
forming \(R^{-1}\).  Define the one-sided resolvents
\begin{equation}
  A_R
  =
  \I-\eta_1RP,
  \qquad
  A_L
  =
  \I-\eta_1PR.
  \label{eq:one-sided-resolvents}
\end{equation}
If \(R\) is invertible, then
\begin{equation}
  R^{-1}-\eta_1P
  =
  R^{-1}A_R
  =
  A_LR^{-1},
  \label{eq:inverse-factorizations}
\end{equation}
and hence
\begin{equation}
  \bigl(R^{-1}-\eta_1P\bigr)^{-1}
  =
  A_R^{-1}R
  =
  RA_L^{-1}.
  \label{eq:resolvent-inverse-equivalence}
\end{equation}
The final resolvent formulas below only require the corresponding one-sided
resolvent to be invertible; no explicit inverse of \(R\) need be formed in
their numerical implementation.
The Sherman--Morrison identity for \(A_R=\I-\eta_1RP\) gives
\begin{equation}
  A_R^{-1}
  =
  \I-\frac{RP}{d},
  \qquad
  A_L^{-1}
  =
  \I-\frac{PR}{d},
  \label{eq:resolvent-SM-expanded}
\end{equation}
where \(d=\Tr(PR)-\eta_1\).  Substitution into
Eqs.~\eqref{eq:LL-invariant}--\eqref{eq:RR-invariant} yields the equivalent
resolvent representation
\begin{align}
  \boxed{
  \Sigma_{LL}'
  =
  \Sigma_{LL}
  +
  \eta_1\Sigma_{LR}P A_R^{-1}\Sigma_{RL}
  }
  &=
  \boxed{
  \Sigma_{LL}
  +
  \eta_1\Sigma_{LR}A_L^{-1}P\Sigma_{RL}
  },
  \label{eq:LL-resolvent}
  \\[5pt]
  \boxed{
  \Sigma_{LR}'
  =
  \Sigma_{LR}Q
  +
  \eta_1\Sigma_{LR}P A_R^{-1}RQ
  }
  &=
  \boxed{
  \Sigma_{LR}A_L^{-1}Q
  },
  \label{eq:LR-resolvent}
  \\[5pt]
  \boxed{
  \Sigma_{RL}'
  =
  Q A_R^{-1}\Sigma_{RL}
  }
  &=
  \boxed{
  QRA_L^{-1}\eta_1P\Sigma_{RL}
  +
  Q\Sigma_{RL}
  },
  \label{eq:RL-resolvent}
  \\[5pt]
  \boxed{
  \Sigma_{RR}'
  =
  \eta_2P
  +
  Q A_R^{-1}RQ
  }
  &=
  \boxed{
  \eta_2P
  +
  QRA_L^{-1}Q
  }.
  \label{eq:RR-resolvent}
\end{align}
For example,
\[
  Q A_R^{-1}\Sigma_{RL}
  =
  Q\Sigma_{RL}
  -
  \frac{QRP\Sigma_{RL}}{d},
  \qquad
  \eta_1\Sigma_{LR}P A_R^{-1}\Sigma_{RL}
  =
  -\frac{\Sigma_{LR}P\Sigma_{RL}}{d},
\]
so the lower-left and upper-left blocks in
Eqs.~\eqref{eq:RL-resolvent} and \eqref{eq:LL-resolvent} reduce exactly to
Eqs.~\eqref{eq:RL-invariant} and \eqref{eq:LL-invariant}.  The remaining
two reductions are identical:
\[
  \eta_1\Sigma_{LR}P A_R^{-1}RQ
  =
  -\frac{\Sigma_{LR}PRQ}{d},
  \qquad
  Q A_R^{-1}RQ
  =
  QRQ-\frac{QRPRQ}{d}.
\]

Equations~\eqref{eq:LL-resolvent}--\eqref{eq:RR-resolvent} are useful
computationally because they replace the explicit scalar division by two
linear solves with the same nearly identity rank-one matrix.  In the
right-resolvent form, solve
\begin{equation}
  A_RX=\Sigma_{RL},
  \qquad
  A_RY=RQ.
  \label{eq:regularized-resolvent-solves}
\end{equation}
Then the update is assembled as
\begin{align}
  \Sigma_{LL}'&=\Sigma_{LL}+\eta_1\Sigma_{LR}PX,
  &
  \Sigma_{LR}'&=\Sigma_{LR}Q+\eta_1\Sigma_{LR}PY,
  \nonumber\\
  \Sigma_{RL}'&=QX,
  &
  \Sigma_{RR}'&=\eta_2P+QY.
  \label{eq:regularized-resolvent-assembly}
\end{align}
The solves in Eq.~\eqref{eq:regularized-resolvent-solves} may be evaluated
with the same Tikhonov-regularized linear solver used for the ordinary
covariance updates.  On well-conditioned nonsingular branches the
unregularized solve gives the exact identities above, while the Tikhonov
fallback supplies a controlled numerical continuation when \(A_R\) is nearly
singular.

\subsection{Off-Diagonal Choi Block}

The upper-right block deserves separate emphasis:
\begin{equation}
  \Sigma_{LR}'
  =
  \Sigma_{LR}Q
  -
  \frac{\Sigma_{LR}PRQ}{d}.
  \label{eq:offdiag-interpretation}
\end{equation}
The first term removes the outgoing column in the measured \(P\)-direction:
\(\Sigma_{LR}'P=0\).  The second term is a finite rank-one correction created
by pre-existing correlations between the selected mode and \(\Ran Q\).  If
\(PRQ=0\), meaning that the selected mode is uncorrelated with its
orthogonal complement in \(R\), then
\begin{equation}
  \Sigma_{LR}'=\Sigma_{LR}Q.
  \label{eq:offdiag-decoupled}
\end{equation}
The analogous statement for the lower-left block follows either directly or
by Hermiticity.

\subsection{Hermiticity and the Singular Branch Condition}

When \(R=R^\dagger\), \(\Sigma_{RL}=\Sigma_{LR}^{\dagger}\), and
\(\eta_{1,2}\) are real, \(d\) is real.  It then follows immediately from
Eq.~\eqref{eq:full-invariant-update} that
\begin{equation}
  \Sigma_{LL}'=(\Sigma_{LL}')^\dagger,
  \qquad
  \Sigma_{RR}'=(\Sigma_{RR}')^\dagger,
  \qquad
  \Sigma_{RL}'=(\Sigma_{LR}')^\dagger.
  \label{eq:hermiticity-preservation}
\end{equation}

If \(d=0\), the rank-one denominator in the regulated contraction vanishes.
The ordinary finite limit in Proposition~\ref{prop:reg-contraction} is then
not defined.  At the operator level this identifies a singular selected
branch: the projected mode has exactly the covariance value at which the
chosen local Kraus update has vanishing or nonregular overlap with the
incoming accumulated tensor.  Such an event requires a separate
zero-overlap or rescaled-limit analysis and is not covered by the formula
above.

\subsection{Preservation of the Choi Involution}

For a pure Gaussian Choi tensor, the covariance is an involution,
\begin{equation}
  \Sigma^2=\I.
  \label{eq:old-choi-involution}
\end{equation}
The corrected regulated update preserves this constraint on every
nonsingular branch.

\begin{proposition}[Preservation of the Choi involution]
  Assume the hypotheses of Proposition~\ref{prop:reg-contraction}, together
  with Eq.~\eqref{eq:old-choi-involution}.  Then, for
  \(d=\Tr(P\Sigma_{RR})-\eta_1\neq0\),
  \begin{equation}
    \left(\Sigma'\right)^2=\I.
    \label{eq:new-choi-involution}
  \end{equation}
\end{proposition}

\begin{proof}
  The assertion is basis independent, so for the proof choose a basis
  adapted to \(P\oplus Q\), and reorder the full covariance as
  \((L\oplus Q)\oplus P\).  Writing
  \begin{equation}
    \Sigma_{LR}=\begin{pmatrix} u & V \end{pmatrix},
    \qquad
    \Sigma_{RR}=
    \begin{pmatrix}
      r & b^\dagger\\
      b & C
    \end{pmatrix},
    \qquad
    r=\Tr(P\Sigma_{RR}),
  \end{equation}
  the unupdated covariance becomes
  \begin{equation}
    \widetilde{\Sigma}
    =
    \begin{pmatrix}
      X & y\\
      y^\dagger & r
    \end{pmatrix},
    \qquad
    X=
    \begin{pmatrix}
      \Sigma_{LL} & V\\
      V^\dagger & C
    \end{pmatrix},
    \qquad
    y=
    \begin{pmatrix}
      u\\ b
    \end{pmatrix}.
    \label{eq:involution-proof-decomposition}
  \end{equation}
  Equation~\eqref{eq:old-choi-involution} is equivalent to
  \begin{align}
    X^2+yy^\dagger &= \I_{L\oplus Q},&
    Xy+r y &=0,&
    y^\dagger y+r^2&=1.
    \label{eq:old-involution-identities}
  \end{align}
  Since \(\widetilde{\Sigma}\) is Hermitian, the second identity also gives
  \(y^\dagger X=-r y^\dagger\).

  The corrected update in Eq.~\eqref{eq:full-invariant-update}, written in
  the same reordered output basis, decouples the newly reset \(P\) mode:
  \begin{equation}
    \widetilde{\Sigma}'
    =
    \begin{pmatrix}
      T & 0\\
      0 & \eta_2
    \end{pmatrix},
    \qquad
    T=X-\frac{yy^\dagger}{d},
    \qquad
    d=r-\eta_1.
    \label{eq:updated-involution-decomposition}
  \end{equation}
  Using Eq.~\eqref{eq:old-involution-identities},
  \begin{align}
    T^2
    &=
    X^2-\frac{Xyy^\dagger+yy^\dagger X}{d}
    +\frac{(yy^\dagger)^2}{d^2}
    \notag\\
    &=
    \I_{L\oplus Q}
    +
    \left(
      -1+\frac{2r}{d}+\frac{1-r^2}{d^2}
    \right)yy^\dagger
    \notag\\
    &=
    \I_{L\oplus Q},
    \label{eq:T-square-identity}
  \end{align}
  because
  \begin{equation}
    -d^2+2rd+1-r^2
    =
    1-\eta_1^2
    =
    0.
  \end{equation}
  Finally, \(\eta_2^2=1\), and therefore
  \(\left(\widetilde{\Sigma}'\right)^2=\I\).  Permuting back to the original
  Choi-leg ordering proves Eq.~\eqref{eq:new-choi-involution}.
\end{proof}

Thus an admissible local branch maps a pure Choi covariance to another pure
Choi covariance: the selected mode is reset to an exact \(\eta_2\)
eigenmode, while the complementary block remains an involution through the
finite Schur-complement update.

\section{Particle--Hole Transfer Complementarity of a Pure Choi Covariance}
\label{sec:bipartite-reconstruction}

The accumulated covariance \(\Sigma\) is itself the covariance of a pure
Gaussian Choi state.  It therefore contains a single-particle transfer
description of the realized non-unitary operator.  The natural language is
not that of a failure of invertibility under projection, but that of
complementary particle and hole amplitudes: an exact projected direction can
be dark for particle propagation while it is a pole of the hole
representation, or vice versa.

In this section the two Choi legs are equal-dimensional one-particle spaces,
\begin{equation}
  \dim\cH_A=\dim\cH_B=N,
  \label{eq:equal-choi-leg-dimensions}
\end{equation}
and the symbols \(A,B,D\) denote Choi-covariance blocks.  They are distinct
from the rank-one measurement projector \(P\) of
Eqs.~\eqref{eq:PQ}--\eqref{eq:full-invariant-update}.

\subsection{Pure covariance identities}

Let \(\ket{\Sigma}\) denote the pure number-conserving Choi state, with
\begin{equation}
  C_{ij}=\bra{\Sigma}c_i^\dagger c_j\ket{\Sigma},
  \qquad
  \Sigma_{ij}=\bra{\Sigma}[c_i^\dagger,c_j]\ket{\Sigma}
  =2C_{ij}-\delta_{ij}.
  \label{eq:bipartite-Gamma-convention}
\end{equation}
The purity condition \(C^2=C\) is equivalent to
\begin{equation}
  \Sigma=\Sigma^\dagger,
  \qquad
  \Sigma^2=\I.
  \label{eq:bipartite-purity}
\end{equation}
After acting on the Choi state, creation and annihilation operators obey
\begin{equation}
  \bm c^\dagger\ket{\Sigma}
  =
  -\Sigma\bm c^\dagger\ket{\Sigma},
  \qquad
  \bm c\,\ket{\Sigma}
  =
  \Sigma^T\bm c\,\ket{\Sigma}.
  \label{eq:global-ket-identities}
\end{equation}
These are state identities and are not operator equalities on the full Fock
space.

\subsection{Consistency of the Choi-state identities}

We now verify explicitly that Eq.~\eqref{eq:global-ket-identities} encodes
the same complex covariance \(\Sigma\) used in
Eq.~\eqref{eq:bipartite-Gamma-convention}.  Sandwiching the annihilation
identity between \(\bra{\Sigma}c_i^\dagger\) and \(\ket{\Sigma}\) gives
\begin{align}
  C_{ij}
  &=
  \bra{\Sigma}c_i^\dagger c_j\ket{\Sigma}
  =
  \sum_k \Sigma_{kj}
  \bra{\Sigma}c_i^\dagger c_k\ket{\Sigma}
  \nonumber\\
  &=
  \sum_k C_{ik}\Sigma_{kj},
  \qquad\Longrightarrow\qquad
  C=C\Sigma .
  \label{eq:annihilation-identity-recovers-occupied-sector}
\end{align}
Because both \(C\) and \(\Sigma\) are Hermitian, taking the adjoint yields
\begin{equation}
  \Sigma C=C.
  \label{eq:occupied-plus-sector}
\end{equation}
Thus the occupied one-particle subspace lies in the \(+1\) eigenspace of
\(\Sigma\).

Similarly, use the creation identity and contract with
\(\bra{\Sigma}c_i\).  Since
\(
  \bra{\Sigma}c_i c_j^\dagger\ket{\Sigma}
  =
  \delta_{ij}-C_{ji}
\),
one obtains
\begin{align}
  \delta_{ij}-C_{ji}
  &=
  -\sum_k\Sigma_{jk}
  \bigl(\delta_{ik}-C_{ki}\bigr).
  \label{eq:creation-contraction-components}
\end{align}
Upon transposing this matrix equation,
\begin{equation}
  \Sigma(\I-C)=-(\I-C).
  \label{eq:empty-minus-sector}
\end{equation}
Hence the empty one-particle subspace lies in the \(-1\) eigenspace of
\(\Sigma\).  Resolving the identity into occupied and empty sectors now
gives
\begin{align}
  \Sigma
  &=
  \Sigma C+\Sigma(\I-C)
  =
  C-(\I-C)
  =
  2C-\I .
  \label{eq:state-identities-recover-covariance}
\end{align}
Therefore
\begin{equation}
  \boxed{
  \bra{\Sigma}[c_i^\dagger,c_j]\ket{\Sigma}
  =
  2\bra{\Sigma}c_i^\dagger c_j\ket{\Sigma}
  -
  \delta_{ij}
  =
  \Sigma_{ij}.
  }
  \label{eq:explicit-complex-covariance-check}
\end{equation}
This closes the consistency loop: the linear relations in
Eq.~\eqref{eq:global-ket-identities} are precisely the state-space
representation of the complex Choi covariance, rather than an additional
transfer assumption.

In the bipartition \(\cH_A\oplus\cH_B\), write
\begin{equation}
  \Sigma
  =
  \begin{pmatrix}
    A & B\\
    B^\dagger & D
  \end{pmatrix},
  \qquad
  A=A^\dagger,
  \qquad
  D=D^\dagger .
  \label{eq:bipartite-blocks}
\end{equation}
Equation~\eqref{eq:bipartite-purity} gives
\begin{align}
  A^2+BB^\dagger&=\I_A,
  &
  AB+BD&=0,
  \label{eq:bipartite-involution-identities}\\
  B^\dagger A+DB^\dagger&=0,
  &
  B^\dagger B+D^2&=\I_B.
  \notag
\end{align}
In particular,
\begin{equation}
  \I_A-A^2=BB^\dagger\geq0.
  \label{eq:A-spectral-bound}
\end{equation}
Since \(A\) is Hermitian, it admits an orthonormal eigenbasis and
\begin{equation}
  \operatorname{spec}(A)\subseteq[-1,1].
  \label{eq:A-spectrum-interval}
\end{equation}

\subsection{Particle and hole transfer matrices}

Taking the \(A\)-leg components of Eq.~\eqref{eq:global-ket-identities}
yields
\begin{align}
  (\I_A+A)\bm c_A^\dagger\ket{\Sigma}
  &=-B\bm c_B^\dagger\ket{\Sigma},
  \label{eq:creation-leg-equation}\\
  (\I_A-A^T)\bm c_A\ket{\Sigma}
  &=B^*\bm c_B\ket{\Sigma}.
  \label{eq:annihilation-leg-equation}
\end{align}
Accordingly, wherever the indicated finite representation exists, define
\begin{align}
  \bm c_A^\dagger\ket{\Sigma}
  &=
  \mathcal T_p\bm c_B^\dagger\ket{\Sigma},
  &
  \mathcal T_p
  &:=
  -(\I_A+A)^{-1}B,
  \label{eq:particle-transfer-definition}\\
  \bm c_A\ket{\Sigma}
  &=
  \mathcal T_h\bm c_B\ket{\Sigma},
  &
  \mathcal T_h
  &:=
  (\I_A-A^T)^{-1}B^* .
  \label{eq:hole-transfer-definition}
\end{align}
Here \(\mathcal T_p\) is the particle, or creation-sector, transfer matrix,
while \(\mathcal T_h\) is the hole, or annihilation-sector, transfer matrix.
Neither definition requires \(A\) itself to have full rank, and a
rank-deficient \(\mathcal T_p\) is entirely regular.  Rather,
\(\mathcal T_p\) develops a pole only at an exact \(A=-\I\) direction,
whereas \(\mathcal T_h\) develops the complementary pole at an exact
\(A=+\I\) direction.  At such endpoints the displayed inverse formulas are
understood by approach from the finite sector.

\begin{proposition}[Particle-transfer representation]
\label{prop:particle-covariance}
On a finite particle-transfer sector, Eq.~\eqref{eq:particle-transfer-definition}
implies
\begin{align}
  A
  &=
  (\I_A-\mathcal T_p\mathcal T_p^\dagger)
  (\I_A+\mathcal T_p\mathcal T_p^\dagger)^{-1},
  \label{eq:particle-A-reconstruction}\\
  B
  &=
  -2(\I_A+\mathcal T_p\mathcal T_p^\dagger)^{-1}\mathcal T_p.
  \label{eq:particle-B-reconstruction}
\end{align}
On the corresponding coupled subspace of \(\cH_B\), one furthermore has
\begin{equation}
  D
  =
  -(\I_B-\mathcal T_p^\dagger\mathcal T_p)
  (\I_B+\mathcal T_p^\dagger\mathcal T_p)^{-1}.
  \label{eq:particle-D-reconstruction}
\end{equation}
\end{proposition}

\begin{proof}
Equation~\eqref{eq:particle-transfer-definition} gives
\begin{equation}
  B=-(\I_A+A)\mathcal T_p.
  \label{eq:B-from-transfer}
\end{equation}
Because \(A=A^\dagger\), substitution into
\(A^2+BB^\dagger=\I_A\) first gives the unsimplified relation
\begin{equation}
  (\I_A+A)\mathcal T_p\mathcal T_p^\dagger(\I_A+A)
  =
  \I_A-A^2 .
  \label{eq:TpTp-congruence}
\end{equation}
On a finite particle-transfer sector, \(\I_A+A\) is invertible.  Multiplying
Eq.~\eqref{eq:TpTp-congruence} on both sides by
\((\I_A+A)^{-1}\) therefore produces the two-inverse expression
\begin{equation}
  \mathcal T_p\mathcal T_p^\dagger
  =
  (\I_A+A)^{-1}
  (\I_A-A^2)
  (\I_A+A)^{-1}.
  \label{eq:TpTp-two-inverses}
\end{equation}
Now \(A\), \(\I_A-A\), \(\I_A+A\), and
\((\I_A+A)^{-1}\) mutually commute, because each is a function of the
Hermitian operator \(A\).  Factoring
\(
  \I_A-A^2=(\I_A-A)(\I_A+A)
\)
and cancelling one factor yields
\begin{equation}
  \mathcal T_p\mathcal T_p^\dagger
  =
  (\I_A-A)(\I_A+A)^{-1}.
  \label{eq:TpTp-from-A}
\end{equation}
The cancellation is unavailable precisely on an \(a=-1\) eigenvector,
where \(\I_A+A\) vanishes and the particle-transfer representation has its
endpoint pole.
Inverting this Cayley relation yields
Eq.~\eqref{eq:particle-A-reconstruction}; inserting the result into
Eq.~\eqref{eq:B-from-transfer} yields
Eq.~\eqref{eq:particle-B-reconstruction}.  Finally, \(AB+BD=0\) fixes
\(D\) on \(\Ran B^\dagger\), giving
Eq.~\eqref{eq:particle-D-reconstruction}.
\end{proof}

It is important that the qualifier in the last statement not be suppressed.
If an exact projection makes \(B\) vanish on a direction, the equation
\(AB+BD=0\) cannot determine the polarization of the decoupled mode on the
opposite Choi leg: there one has only \(D^2=\I_B\).  Thus
Eqs.~\eqref{eq:particle-A-reconstruction} and
\eqref{eq:particle-B-reconstruction} determine particle transfer data,
whereas a fully projected endpoint retains additional output-polarization
information in \(D\).  This is precisely the information carried by the
\(\eta_2P\) reset term in Eq.~\eqref{eq:full-invariant-update}.

\subsection{Polar form and particle--hole duality}

Take the left polar decomposition
\begin{equation}
  \mathcal T_p=HU,
  \qquad
  H=(\mathcal T_p\mathcal T_p^\dagger)^{1/2}\geq0.
  \label{eq:left-polar-transfer}
\end{equation}
Because the Choi legs have equal dimension, the partial isometry can be
extended to a unitary \(U:\cH_B\to\cH_A\).  On the particle-transfer
sector,
\begin{equation}
  A=q(H):=\frac{\I-H^2}{\I+H^2},
  \qquad
  B=-r(H)U,
  \qquad
  r(H):=\frac{2H}{\I+H^2}.
  \label{eq:qr-functions}
\end{equation}
The functions \(q(H)\) and \(r(H)\) commute and obey
\begin{equation}
  q(H)^2+r(H)^2=\I_A,
  \label{eq:qr-circle}
\end{equation}
making purity transparent on the coupled sector.

For strictly positive finite singular values of \(H\), substitution into
Eq.~\eqref{eq:hole-transfer-definition} gives
\begin{equation}
  \mathcal T_h
  =
  -\left(H^T\right)^{-1}U^*
  =
  -H^{-T}U^* .
  \label{eq:particle-hole-transfer-duality}
\end{equation}
Thus particle and hole amplitudes are dual descriptions of the same pure
Choi correlation:
\begin{equation}
  s_j(\mathcal T_h)
  =
  s_j(\mathcal T_p)^{-1}
  \qquad
  \text{for }0<s_j(\mathcal T_p)<\infty .
  \label{eq:reciprocal-singular-values}
\end{equation}
This is particle--hole complementarity in transfer space, rather than an
alternative convention appended to the dynamics.

\subsection{Endpoint sectors and Lyapunov exponents}

Let
\begin{equation}
  A u_j=a_j u_j,
  \qquad
  a_j\in[-1,1].
  \label{eq:A-eigenproblem}
\end{equation}
Equation~\eqref{eq:TpTp-from-A} gives, on the finite sector,
\begin{equation}
  s_{p,j}^2
  =
  \frac{1-a_j}{1+a_j},
  \qquad
  s_{h,j}^2
  =
  \frac{1+a_j}{1-a_j}.
  \label{eq:particle-hole-singular-from-A}
\end{equation}
Consequently the endpoint structure is
\begin{align}
  a_j=+1:
  &\qquad
  s_{p,j}=0,
  \qquad
  s_{h,j}=\infty,
  \label{eq:plus-endpoint}\\
  -1<a_j<+1:
  &\qquad
  0<s_{p,j}<\infty,
  \qquad
  s_{h,j}=s_{p,j}^{-1},
  \label{eq:interior-sector}\\
  a_j=-1:
  &\qquad
  s_{p,j}=\infty,
  \qquad
  s_{h,j}=0.
  \label{eq:minus-endpoint}
\end{align}
The infinite values in Eqs.~\eqref{eq:plus-endpoint} and
\eqref{eq:minus-endpoint} are limiting statements: at the exact endpoint
the relevant inverse in Eq.~\eqref{eq:particle-transfer-definition} or
\eqref{eq:hole-transfer-definition} is singular.  Exact projective
operations polarize selected one-particle directions into these extremal
covariance sectors.  A particle-dark direction is a pole in the
complementary hole description, and a hole-dark direction is a pole in the
particle description.  This is not a loss of purity or a failure of the
Choi covariance; it is the sharp transfer signature of an exact Gaussian
projection.

For an accumulated Choi tensor after \(t>0\) completed cycles, the finite
particle and hole Lyapunov exponents are therefore
\begin{align}
  \lambda_{p,j}(t)
  &=
  \frac{1}{t}\log s_{p,j}(t)
  =
  \frac{1}{2t}
  \log\!\left(\frac{1-a_j(t)}{1+a_j(t)}\right)
  =
  -\frac{1}{t}\operatorname{artanh}a_j(t),
  \label{eq:particle-lyapunov-from-A}\\
  \lambda_{h,j}(t)
  &=
  \frac{1}{t}\log s_{h,j}(t)
  =
  -\lambda_{p,j}(t).
  \label{eq:hole-lyapunov-from-A}
\end{align}
The endpoint sectors contribute formal \(-\infty\) or \(+\infty\)
exponents in the appropriate particle or hole representation.  A supported
numerical observable may omit these endpoint sectors and record their
multiplicities separately, but it must identify the transfer convention to
which its retained exponents refer.

\subsection{Particle-transfer Lyapunov gap}

The transfer-space gap relevant for the slowest finite particle direction is
the finite-time Lyapunov gap
\begin{equation}
  \Delta_\lambda(t)
  \equiv
  \min_{j\ {\rm finite}}\left|\lambda_{p,j}(t)\right|
  =
  \frac{1}{t}
  \min_{j\ {\rm finite}}
  \left|\log s_j\!\left(\mathcal T_p(t)\right)\right| .
  \label{eq:particle-lyapunov-gap-definition}
\end{equation}
By Eq.~\eqref{eq:particle-lyapunov-from-A}, it is obtained directly from the
left Choi block:
\begin{equation}
  \boxed{
  \Delta_\lambda(t)
  =
  \frac{1}{t}
  \min_{j\ {\rm finite}}
  \left|
    \operatorname{artanh} a_j(t)
  \right|,
  \qquad
  a_j(t)\in\operatorname{spec}\Sigma_{LL}(t).
  }
  \label{eq:particle-gap-from-A}
\end{equation}
On the open interval \((-1,1)\), the function
\(\left|\operatorname{artanh}a\right|\) is strictly increasing with
\(\lvert a\rvert\).  Thus the finite gap is controlled by eigenvalues of
\(A(t)=\Sigma_{LL}(t)\) closest to zero, rather than by its polarized
endpoint sectors.  The latter have \(a_j=\pm1\) and contribute formal
infinite particle exponents; they are excluded from the finite minimum in
Eq.~\eqref{eq:particle-gap-from-A} while their multiplicities remain a
separate diagnostic of exact projection.

Equivalently, if \(A u_j=a_j u_j\), then
Eq.~\eqref{eq:TpTp-from-A} gives
\begin{equation}
  \mathcal T_p\mathcal T_p^\dagger u_j
  =
  \frac{1-a_j}{1+a_j}u_j
  =
  s_{p,j}^2 u_j.
  \label{eq:near-gap-transfer-eigenstates}
\end{equation}
The \(A\)-eigenstates closest to \(a_j=0\) are therefore precisely the
eigenstates of \(\mathcal T_p\mathcal T_p^\dagger\) whose eigenvalues are
closest to unity.  These are the near-gap transfer states.

It is important not to replace this criterion by
\(\sigma_{\min}(\I-\mathcal T_p)\).  The latter probes whether the operator
\(\mathcal T_p\) has an approximately fixed vector, whereas
Eq.~\eqref{eq:particle-lyapunov-gap-definition} probes whether a singular
value of \(\mathcal T_p\) lies near unity.  For a non-normal transfer
operator these questions are inequivalent.  Their identification would
require additional normality and alignment assumptions that do not follow
from purity of the Choi covariance.

\subsection{Information content of the Choi blocks}

For the particle Lyapunov spectrum, the upper-left Choi block is already
complete data.  Indeed, Eq.~\eqref{eq:TpTp-from-A} states that
\begin{equation}
  \mathcal T_p\mathcal T_p^\dagger
  =
  (\I_A-A)(\I_A+A)^{-1}.
  \label{eq:TpTpd-eigenstates-from-A}
\end{equation}
Consequently the eigenvalues of \(A=\Sigma_{LL}\) determine all particle
singular values, while its eigenvectors are precisely the eigenstates of
the positive transfer operator \(\mathcal T_p\mathcal T_p^\dagger\)
(equivalently, the left singular vectors of \(\mathcal T_p\)).  This
includes the regular zero-eigenvalue sector at \(a=+1\).  Thus neither an
explicit construction of
\(\mathcal T_p\) nor a spectral decomposition of \(B=\Sigma_{LR}\) is
required in order to compute \(\lambda_{p,j}(t)\) and the corresponding
\(\mathcal T_p\mathcal T_p^\dagger\) eigenstates.

For a complete particle spectrum one may diagonalize \(A=\Sigma_{LL}\).
For the gap alone, however, a complete intermediate spectral resolution is
unnecessary.  Since the near-gap states minimize \(\lvert a_j\rvert\), they
are the lowest-eigenvalue states of the positive operator
\begin{equation}
  A^2=\Sigma_{LL}^2 .
  \label{eq:matrix-free-gap-target}
\end{equation}
An iterative Hermitian eigensolver can access this sector using only the
matrix-vector operation
\begin{equation}
  v\longmapsto A(Av),
  \label{eq:matrix-free-gap-action}
\end{equation}
without materializing \(A^2\), \(\mathcal T_p\), or
\(\mathcal T_p\mathcal T_p^\dagger\).  Because \(A^2\) does not distinguish
the signs of paired eigenvalues \(+a\) and \(-a\), the signed particle
exponents must be recovered by diagonalizing the projected Hermitian
operator \(X^\dagger A X\) inside the converged near-zero subspace \(X\).
An exact diagonalization of \(A\) at the final observation time then provides
the complete final spectrum and an independent validation of the iterative
near-gap extraction.

The off-diagonal block has different, equally physical information content.
In the polar representation of Eq.~\eqref{eq:qr-functions}, \(A\) fixes
the positive factor \(H\), whereas
\begin{equation}
  B=-r(H)U
  \label{eq:B-carries-polar-isometry}
\end{equation}
fixes the partial-isometry, or angular, data \(U\) on the coupled support.
Therefore \(B\) is required to reconstruct the full complex transfer matrix,
its input-leg spectral data, or the mapping between input-leg and output-leg
orbitals.  The distinction is operationally useful: a calculation whose
target is only the Lyapunov spectrum and the output-leg eigenstates of
\(\mathcal T_p\mathcal T_p^\dagger\) may diagonalize \(\Sigma_{LL}\)
alone, whereas an analysis of the transported wavefunctions themselves must
retain \(\Sigma_{LR}\).

\subsection{Regularity of a zero particle-transfer sector}

Rank deficiency of \(\mathcal T_p\) itself is regular.  Let \(\Pi_0\) be the
orthogonal projector onto \(\ker H\).  Then
\begin{equation}
  q(H)\Pi_0=\Pi_0,
  \qquad
  r(H)\Pi_0=0,
  \label{eq:kernel-qr}
\end{equation}
so \(H=0\) is precisely the \(a=+1\), particle-dark endpoint.  If desired,
a full-rank family approaching this endpoint may be introduced by
\begin{equation}
  H_\eps=H+\eps\Pi_0,
  \qquad
  \eps>0.
  \label{eq:H-kernel-regulator}
\end{equation}
Restricted to \(\Ran\Pi_0\),
\begin{align}
  q(H_\eps)\Pi_0
  &=
  \left(1-2\eps^2+O(\eps^4)\right)\Pi_0,
  \label{eq:q-kernel-regulated}\\
  r(H_\eps)\Pi_0
  &=
  \left(2\eps+O(\eps^3)\right)\Pi_0.
  \label{eq:r-kernel-regulated}
\end{align}
The particle amplitude therefore vanishes continuously while its hole
partner diverges in accordance with
Eq.~\eqref{eq:reciprocal-singular-values}.

\begin{remark}
The matrices \(\mathcal T_p\) and \(\mathcal T_h\) above are complex,
number-conserving particle and hole transfer matrices extracted from
\(\Sigma\).  A direct complex-basis GPU observable associated with
\path{colab_regularized_choi_transfer_matrix/src/choi_transfer_observables_gpu.py}
uses the particle-spectrum consequence of
Eq.~\eqref{eq:TpTpd-eigenstates-from-A} in the complex basis.  For
time-resolved gap extraction it may target the near-zero spectrum of
\(A^2\) through Eq.~\eqref{eq:matrix-free-gap-action}, while an exact
final-time diagonalization of \(A\) supplies the full final particle
spectrum and selected near-gap eigenstates of
\(\mathcal T_p\mathcal T_p^\dagger\).  No Majorana covariance,
pseudoinverse, or explicit stored \(\mathcal T_p\) is required.
The complementary finite hole exponents are recovered by
Eq.~\eqref{eq:hole-lyapunov-from-A}.
\end{remark}

\section{Relation to the Repository Dynamics}

The cumulative-Choi construction developed here is useful when one intends to
retain a composed non-unitary operator and, for example, extract a transfer
spectrum from its Choi blocks.  Its complex blocks encode the complementary
particle and hole transfer data described in
Sec.~\ref{sec:bipartite-reconstruction}.  The dedicated analysis bundle,
\path{colab_regularized_choi_transfer_matrix/}, propagates the corrected
cumulative Choi blocks.  The intended finite-size gap campaign uses
\begin{equation}
  N_x=20,
  \qquad
  N_y\in\{30,40,50\},
  \qquad
  T=2N_y ,
  \label{eq:intended-gap-campaign}
\end{equation}
and tracks
\begin{equation}
  \Delta_\lambda(t)=\min_j\left|\lambda_{p,j}(t)\right|
  \label{eq:repository-gap-observable}
\end{equation}
at each completed cycle.  In the domain-wall truncated geometry the
observable is evaluated on the active slab-only Choi basis; in the
homogeneous \(DW=\mathrm{false}\) geometry it is evaluated on the complete
top-layer Choi basis.  The numerical strategy is to extract the near-gap
sector through the matrix-free \(A^2\) procedure above throughout the
trajectory, and at \(t=T\) to retain an exact particle spectrum and selected
near-gap eigenstates of \(\mathcal T_p\mathcal T_p^\dagger\).  This describes
the intended direct complex-basis observable and does not assert that these
finite-size campaign products have already been generated.

The canonical Markov simulations in this repository take a different route.
The CPU and GPU implementations of
\path{classA_U1FGTN.run_markov_circuit(...)} and
\path{classA_U1FGTN_gpu.run_markov_circuit(...)} propagate
the normalized top-layer covariance \(G_t\), not \(\Sigma_t\).  For finite
Wannier support, their rank-one measurement maps are evaluated by solving on
the compact support block of the local mode.  Those support-reduced formulas
are derived in
\path{src/fgtn/support_reduced_markov_update_note.tex}.  The stochastic
covariance process and its tangent cocycle are organized in
\path{markov_transfer_operators/markov_transfer_operators_note.tex}.

Accordingly, distinct constructions should not be conflated:
\begin{enumerate}
  \item The present \(\eps\)-prescription is an analytic definition of a
  singular cumulative-Choi contraction before taking \(\eps\to0\).
  \item Endpoint poles of \(\mathcal T_p\) or \(\mathcal T_h\) express
  particle--hole complementarity under an exact projective update.
  \item The direct complex particle-spectrum observable extracts only
  \(\mathcal T_p\mathcal T_p^\dagger\) data; reconstructing the full
  transfer matrix or its input-leg spectral data requires the additional
  angular information in \(\Sigma_{LR}\).
\end{enumerate}
Only the first is required to define the regulated contraction itself.

\section{Numerical Consistency Check}

The basis-independent result can be verified without using any circuit
simulation data.  For each test one generates Hermitian
\(\Sigma_{LL}\) and \(R=\Sigma_{RR}\), an arbitrary \(\Sigma_{LR}\) with
\(\Sigma_{RL}=\Sigma_{LR}^\dagger\), a rank-one projector \(P\), and each
of the four sign pairs \((\eta_1,\eta_2)\).  Cases with \(d\) close to zero
are excluded.  The directly regulated contraction obtained at finite
\(\eps\) from Eqs.~\eqref{eq:singular-composition} and
\eqref{eq:regulator} is then compared to
Eq.~\eqref{eq:full-invariant-update}.

For representative random matrices with \(N=3,5,7\), the largest Frobenius
discrepancy over the four sign pairs was
\begin{table}[h]
  \caption{\label{tab:numerical-check}
  Maximum full-Choi-block discrepancy between the regulated finite-\(\eps\)
  contraction and the basis-independent \(\eps\to0\) formula for a fixed
  random test ensemble.  The linear decrease with \(\eps\) is consistent
  with the first omitted term in the regulated expansion.}
  \begin{ruledtabular}
  \begin{tabular}{cccc}
  Dimension \(N\) & \(\eps=10^{-4}\) & \(\eps=10^{-6}\) & \(\eps=10^{-8}\)\\
  \colrule
  \(3\) & \(1.9\times10^{-3}\) & \(1.9\times10^{-5}\) & \(1.9\times10^{-7}\)\\
  \(5\) & \(9.0\times10^{-3}\) & \(9.0\times10^{-5}\) & \(9.0\times10^{-7}\)\\
  \(7\) & \(2.6\times10^{-2}\) & \(2.6\times10^{-4}\) & \(2.6\times10^{-6}\)
  \end{tabular}
  \end{ruledtabular}
\end{table}
Hermiticity of the output blocks is also satisfied to floating-point
precision, in agreement with Eq.~\eqref{eq:hermiticity-preservation}.

The involution statement was tested separately by generating random
Hermitian involutions \(\Sigma=U\operatorname{diag}(\pm1)U^\dagger\), random rank-one
projectors \(P\), and applying all four sign branches while excluding
\(\lvert d\rvert<10^{-8}\).  For \(N=3,5,8\), with \(100\) samples per dimension
and \(1200\) nonsingular branch updates in total, the worst residuals were
\begin{equation}
  \max \left\lVert\left(\Sigma'\right)^2-\I\right\rVert_{F}
  =
  7.317\times10^{-15},
  \qquad
  \max \left\lVert\Sigma'-\left(\Sigma'\right)^\dagger\right\rVert_{F}
  =
  8.689\times10^{-16}.
\end{equation}
These errors are at floating-point roundoff and provide an independent
check of the algebraic preservation proof.

\section{Conclusion}

The local adaptive Kraus branches of the class-A fermionic circuit admit a
compact two-sign Choi representation, but their rank-one projector structure
makes naive full-tensor composition singular.  The regulated contraction
\(K_{LL}^{(\eps)}=\eta_1P+\eps Q\) supplies a controlled limit and exposes a
simple basis-independent answer.  The selected mode is reset by
\(\eta_2P\); all other modes are updated by an embedded Schur complement; and
the inter-leg Choi correlations receive a finite correction determined by the
correlation of the selected mode with its complement.  Crucially, the
lower-right reduced Schur complement must be embedded into \(\Ran Q\) before
it is combined with the full-space \(P\)-sector contribution.  This resolves
the dimensional ambiguity in the handwritten expression and gives a complete
operator-valued update suitable for cumulative Choi transfer analyses.
When the input covariance is pure, \(\Sigma^2=\I\), the corrected update
also preserves the involution exactly on every nonsingular branch.  Finally,
the bipartite Choi covariance resolves the transfer content into conjugate
particle and hole descriptions.  A rank-deficient particle transfer is a
regular, particle-dark sector; its complementary hole representation has an
endpoint pole, and the opposite statement holds at the other extremal
covariance polarization.  Any supported numerical spectrum must therefore
state its basis and whether it retains particle, hole, or Majorana transfer
data.

\end{document}
